\documentclass[11pt]{article}
\usepackage[a4paper,margin=27mm]{geometry}
\usepackage{amsmath,amssymb,amsthm,booktabs,array}
\usepackage[T1]{fontenc}
\usepackage{lmodern}
\usepackage[hidelinks]{hyperref}
\newtheorem{theorem}{Theorem}
\newtheorem{lemma}[theorem]{Lemma}
\newtheorem{corollary}[theorem]{Corollary}
\newcommand{\F}{F}
\title{A rational counterexample to Karp's fractional-degree\newline Legendre Tur\'an conjecture}
\author{Haobo Ma\\\small ChronoAI Pte Ltd, Singapore
\and Wenlin Zhang\\\small National University of Singapore\\\small ChronoAI Pte Ltd, Singapore}
\date{November 15, 2026}
\begin{document}
\maketitle
\begin{abstract}
Karp's Conjecture 4 in his 2011 SF21 presentation
\cite[physical PDF p.~52, repeated on pp.~53--54]{karp} fails at
$\mu=3/2$, $\delta=63/64$, and $y=4095/4096$:
exact rational series bounds give a Legendre Tur\'anian at most
$-97/5000$, with all three degrees positive.
The standard zero-balanced hypergeometric expansion also gives
counterexamples at every positive noninteger central degree, for shifts
sufficiently close to one. For each fixed admissible degree and shift,
the Tur\'anian is positive near both ends of $0<y<1$; hence the explicit
counterexample lies between at least two interior zeros.
\end{abstract}

\section{The question and the counterexample}
For real $\nu$ and $0<y<1$, set
\begin{equation}\label{defF}
 \F_\nu(y)={}_2F_1(-\nu,1+\nu;1;y)
 =\sum_{n=0}^{\infty}\frac{(-\nu)_n(1+\nu)_n}{(n!)^2}y^n,
\end{equation}
where $(a)_0=1$ and $(a)_n=a(a+1)\cdots(a+n-1)$ for $n\geq1$.
This is the Ferrers function $\mathsf P_\nu(1-2y)$ in the standard
normalization \cite[\S14.3(i)]{dlmflegendre}; the series convention is
\cite[\S15.2(i)]{dlmfhypergeometric}. Define
\[
 T_{\mu,\delta}(y)=\F_\mu(y)^2-
 \F_{\mu+\delta}(y)\F_{\mu-\delta}(y).
\]
Following the classical integer-degree, unit-shift inequality, Karp
\cite[p.~52, Conjecture~4]{karp} states:
\begin{quote}
``Tur\'an inequality is true for all $\mu>0$ and $0<\delta<1$.''
\end{quote}
Here the displayed inequality on that page is $T_{\mu,\delta}(y)>0$
for $0<y<1$.
The displayed conjecture has no sign restriction on the function values.
Direct series enclosures handle the sign changes that prevent an argument
using logarithms of positive values.

\begin{theorem}\label{thm:main}
Let $y_*=4095/4096$. Then
\begin{align*}
 -\frac{23}{10}&\leq \F_{33/64}(y_*)\leq-\frac{227}{100},\\
 \frac{181}{100}&\leq \F_{3/2}(y_*)\leq\frac{184}{100},\\
 -\frac{153}{100}&\leq \F_{159/64}(y_*)\leq-\frac{150}{100}.
\end{align*}
Consequently,
\begin{equation}\label{neg}
 T_{3/2,\,63/64}(y_*)\leq
 \left(\frac{184}{100}\right)^2-
 \frac{227}{100}\frac{150}{100}
 =-\frac{97}{5000}<0.
\end{equation}
Thus both the strict and non-strict forms of the conjectured comparison
fail, even with $\mu-\delta>0$.
\end{theorem}
The non-strict version is not a separate assertion in the source; it is
a weaker comparison that fails at the same point.

\section{Exact enclosures and proof}
Write $u_n=(-\nu)_n(1+\nu)_ny^n/(n!)^2$. Multiplication of rising
factorials gives, including cases where a term vanishes,
\begin{equation}\label{recurrence}
 u_0=1,\qquad
 u_{n+1}=u_n\,y\frac{(n-\nu)(n+\nu+1)}{(n+1)^2}.
\end{equation}
If $n\geq\nu\geq0$, the last rational factor lies in $[0,1]$:
its numerator is nonnegative and equals
$(n+1)^2-(n+1)-\nu(\nu+1)$. In particular,
$|u_{n+1}|\leq y|u_n|$.

\begin{lemma}\label{lemma:tail}
If $N\geq\nu\geq0$, $0<y<1$, and $S_N=\sum_{n=0}^N u_n$, then
\[
 |\F_\nu(y)-S_N|\leq |u_N|\frac{y}{1-y}.
\]
\end{lemma}
\begin{proof}
Iteration of \eqref{recurrence} gives $|u_{N+j}|\leq |u_N|y^j$
for $j\geq1$. Summing this geometric majorant proves absolute convergence
of the tail and the stated bound.
\end{proof}

Here is a fully rational construction of the required finite bounds.
Let $\nu=p/q$, $y=r/s$, with $q,s>0$, and fix $M=10^{12}$.
Set $L_0=U_0=A_0=B_0=M$, and, for $0\leq n<N$, define
\begin{align*}
 a_n&=(nq-p)((n+1)q+p)r,&
 b_n&=q^2(n+1)^2s,\\
 L_{n+1}&=\left\lfloor\frac{\min(a_nL_n,a_nU_n)}{b_n}\right\rfloor,&
 U_{n+1}&=\left\lceil\frac{\max(a_nL_n,a_nU_n)}{b_n}\right\rceil,\\
 A_{n+1}&=A_n+L_{n+1},& B_{n+1}&=B_n+U_{n+1}.
\end{align*}
Because $b_n>0$, induction using \eqref{recurrence} yields
\begin{equation}\label{invariants}
 L_n\leq Mu_n\leq U_n,\qquad A_n\leq MS_n\leq B_n.
\end{equation}
The minimum and maximum are essential when $a_n$ is negative.
For $N=10\,000$ and $y=y_*$, the integer recurrences give the following
values; \texttt{check\_bounds.py} evaluates exactly these recurrences.

\begin{center}
\small
\begin{tabular}{@{}rrrrr@{}}
\toprule
$\nu$ & $L_N$ & $U_N$ & $A_N$ & $B_N$\\
\midrule
$33/64$ & $-2767853$ & $-2765298$ & $-2284872334619$ & $-2284856488431$\\
$3/2$ & $2769448$ & $2772012$ & $1821604319336$ & $1821620117692$\\
$159/64$ & $-2770040$ & $-2767484$ & $-1516986390784$ & $-1516970637096$\\
\bottomrule
\end{tabular}
\end{center}
With $E=\max(|L_N|,|U_N|)$, Lemma~\ref{lemma:tail},
\eqref{invariants}, and $y_* /(1-y_*)=4095$ give
\begin{equation}\label{enclosure}
 \frac{A_N-4095E}{10^{12}}\leq\F_\nu(y_*)\leq
 \frac{B_N+4095E}{10^{12}}.
\end{equation}
Substitution and integer comparisons give the three bounds of
Theorem~\ref{thm:main}. The central interval is positive, so its square is
at most $(184/100)^2$. The two other intervals are negative, so their
product is at least $(227/100)(150/100)$. This proves \eqref{neg}.
Finally, $0<63/64<1$, $0<y_*<1$, and
$3/2-63/64=33/64>0$, verifying every required parameter restriction.

The file \texttt{check\_bounds.py} evaluates the displayed integer
recurrences and compares the enclosures using exact rational arithmetic;
\texttt{bounds.json} gives their values.

\section{Endpoint behavior and further consequences}
\begin{corollary}\label{cor:open}
There is an open neighborhood of $(3/2,63/64,4095/4096)$ on which
$\mu>\delta>0$, $\delta<1$, $0<y<1$, and $T_{\mu,\delta}(y)<0$.
\end{corollary}
\begin{proof}
Take a compact neighborhood on which all three degrees belong to a
compact subinterval of $(0,\infty)$ and $y\leq q_0<1$.
Once $n$ exceeds every degree in that interval,
\eqref{recurrence} bounds successive absolute terms by $q_0$.
The first finitely many terms are continuous in degree and argument
and uniformly bounded on that compact set. The resulting uniform
geometric majorant proves continuity of each series sum, and hence
of $T$. Strict negativity in \eqref{neg} and the strict parameter
inequalities persist on a smaller open neighborhood.
\end{proof}

\begin{corollary}\label{cor:local}
For every fixed $\mu>0$ and $0<\delta<1$, there is
$\varepsilon=\varepsilon(\mu,\delta)>0$ such that
$T_{\mu,\delta}(y)>0$ for $0<y<\varepsilon$.
\end{corollary}
\begin{proof}
The convergent series at zero gives
$\F_\nu(y)=1-\nu(\nu+1)y+O(y^2)$ for fixed real $\nu$.
Substitution shows
\[
 T_{\mu,\delta}(y)=2\delta^2 y+O(y^2).
\]
The leading coefficient is positive. Dividing by $y>0$ and letting
$y\downarrow0$ proves the result.
\end{proof}
There is also positivity sufficiently close to the other endpoint.
Write $t=1-y$ and $L=-\log t$. For real $\nu$ put
$a_\nu=-\sin(\pi\nu)/\pi$, and for $\nu\notin\mathbb Z$ put
\[
 h_\nu=2\psi(1)-\psi(-\nu)-\psi(1+\nu),
\]
where $\psi=\Gamma'/\Gamma$. The zero-balanced expansion
\cite[Eq.~15.8.10, $m=0$]{dlmfcontinuation} and the reflection formula
\cite[Eq.~5.5.3]{dlmfgamma} give, for fixed noninteger $\nu$,
\begin{equation}\label{eq:asymptotic}
 \F_\nu(y)=a_\nu(L+h_\nu)+O(tL),\qquad t\downarrow0.
\end{equation}
Indeed, the $k=0$ term of that expansion is
$[L+2\psi(1)-\psi(-\nu)-\psi(1+\nu)]/
[\Gamma(-\nu)\Gamma(1+\nu)]$, and
$1/[\Gamma(-\nu)\Gamma(1+\nu)]=-\sin(\pi\nu)/\pi$.
The remaining convergent power series, including its logarithmic factor,
is $O(tL)$: both its analytic coefficient series have zero constant term.
No gamma or digamma pole occurs for these noninteger parameters.
For integer $\nu$, symmetry $\F_\nu=\F_{-\nu-1}$ reduces to a
nonnegative integer degree, whose series terminates. Consequently,
for every fixed real $\nu$, $\F_\nu(y)=a_\nu L+O(1)$.

\begin{corollary}\label{cor:endpoint}
For every fixed $\mu>0$ and $0<\delta<1$,
\begin{equation}\label{eq:endpoint}
 T_{\mu,\delta}(y)=\frac{\sin^2(\pi\delta)}{\pi^2}
                 \log^2(1-y)+O(|\log(1-y)|),\qquad y\uparrow1.
\end{equation}
In particular, $T_{\mu,\delta}$ is positive sufficiently near $y=1$.
\end{corollary}
\begin{proof}
Multiply the three expansions $\F_\nu=a_\nu L+O(1)$ and use
\[
 \sin^2(\pi\mu)-\sin(\pi(\mu+\delta))\sin(\pi(\mu-\delta))
 =\sin^2(\pi\delta)>0.
\]
The error products are $O(L)$, and the leading coefficient is positive.
\end{proof}

\begin{corollary}\label{cor:family}
For each positive noninteger $\mu$, there exist $y_\mu\in(0,1)$
and $\eta_\mu\in(0,1)$ such that
$T_{\mu,\delta}(y_\mu)<0$ for every $1-\eta_\mu<\delta<1$.
If $\mu>1$, all three degrees are positive.
\end{corollary}
\begin{proof}
First allow $\delta=1$. By the digamma recurrence
$\psi(z+1)=\psi(z)+1/z$ \cite[Eq.~5.5.2]{dlmfgamma},
\[
 h_{\mu+1}=h_\mu-\frac{2}{\mu+1},\qquad
 h_{\mu-1}=h_\mu+\frac{2}{\mu},\qquad
 a_{\mu+1}=a_{\mu-1}=-a_\mu.
\]
Substituting \eqref{eq:asymptotic}, with $H=L+h_\mu$, gives
\begin{align*}
 T_{\mu,1}(y)
 &=a_\mu^2\left[H^2-
       \left(H-\frac{2}{\mu+1}\right)
       \left(H+\frac{2}{\mu}\right)\right]+O(tL^2)\\
 &=-\frac{2a_\mu^2}{\mu(\mu+1)}(L+h_\mu-2)+o(1).
\end{align*}
Since $a_\mu\ne0$, this tends to $-\infty$.
Choose a fixed $y_\mu<1$ where it is negative.
For this fixed argument, the hypergeometric series is continuous in
degree: on any compact degree interval, the rational factor in
\eqref{recurrence} tends uniformly to one, so the absolute term ratio
is eventually at most some $\rho<1$. This proves local uniform
convergence and continuity, including at terminating degrees.
Thus negativity persists for all shifts sufficiently close to one from
below. Finally, $\mu-\delta>\mu-1>0$ when $\mu>1$.
\end{proof}
This argument also shows that no endpoint interval of positivity in
Corollary~\ref{cor:endpoint} can be uniform over all $0<\delta<1$
at a fixed positive noninteger $\mu$: choose $y_\mu$ arbitrarily close
to one before applying continuity in the shift.

\begin{corollary}\label{cor:zeros}
The function $T_{3/2,63/64}$ has at least two distinct zeros in $(0,1)$,
one on each side of $4095/4096$.
\end{corollary}
\begin{proof}
Corollaries~\ref{cor:local} and~\ref{cor:endpoint} give positive values
on both sides of the negative value in Theorem~\ref{thm:main}.
Continuity and the intermediate value theorem give the two zeros.
\end{proof}

\section{Relation to the posing source}
The endpoint statements and the family in
Corollary~\ref{cor:family} are not separate conjectures in the slides.
They explain why a fixed-parameter endpoint test does not imply the
global comparison. As a direct boundary calculation,
$\F_0=1$, $\F_1=1-2y$, and $\F_2=1-6y+6y^2$, whence
$T_{1,1}(y)=2y(1-y)>0$ on $(0,1)$.
This is consistent with the classical integer-degree, unit-shift
comparison recalled on p.~52 of \cite{karp}.

The polynomial siblings on p.~31 have an elementary equality boundary.
For $n\geq3$, $\alpha,\beta>0$, and positive $f_0,\ldots,f_n$, write
their polynomial as
\[
 P_{\alpha,\beta}^{n}(x)=\sum_{k=0}^n\binom nk f_kf_{n-k}
 [(x+\alpha)_k(x+\beta)_{n-k}-(x+\alpha+\beta)_k(x)_{n-k}],
\]
take a geometric sequence $f_k=cr^k$ with $c,r>0$.
Both sums then equal $c^2r^n(2x+\alpha+\beta)_n$ by the
rising-factorial Vandermonde identity, so $P_{\alpha,\beta}^{n}\equiv0$.
This equality case was already recorded by Karp
\cite[Lemma~3]{karppolynomial} for $f_k=1$; the geometric case follows
by the constant factor $c^2r^n$.
Thus a positive-coefficient assertion under non-strict log-concavity
must exclude this boundary. The later formulations in
\cite[p.~2, Conjectures~1--2]{karppolynomial} assume strict
log-concavity. For $\alpha=\beta=1$, Karp's Theorem~1 there proves
positive coefficients of $P_{1,1}^{n}(x-1)$, and hence of
$P_{1,1}^{n}(x)$ by translating by one. This is a prior polynomial
result, not a consequence of the Legendre counterexample.

\begin{center}
\small
\begin{tabular}{@{}p{.44\linewidth}p{.49\linewidth}@{}}
\toprule
Source item or neighboring statement & Coverage\\
\midrule
Karp, Conjecture 4, pp.~52--54 & Refuted: Theorem~\ref{thm:main}, even on $\mu>\delta$.\\
Non-strict comparison & Not asserted in the source; fails at the same point.\\
Persistence near the witness & Not posed; proved in Corollary~\ref{cor:open}.\\
Positivity near $y=0$ and $y=1$ & Not posed; proved for each fixed pair in Corollaries~\ref{cor:local} and~\ref{cor:endpoint}.\\
Every positive noninteger central degree & Not posed; a negative interior value for shifts near one is proved in Corollary~\ref{cor:family}.\\
Two interior zeros at the rational parameters & Not posed; proved in Corollary~\ref{cor:zeros}.\\
Integer-degree, unit-shift comparison, p.~52 & Prior result recalled by Karp; the case $\mu=1$ is calculated above.\\
Geometric equality case in the polynomial siblings & Reproved; prior proof \cite[Lemma~3]{karppolynomial} for the constant sequence, extended by scaling above.\\
Conjecture 1, p.~31 (positive polynomial coefficients) & Non-strict reading fails for geometric sequences. Strict $\alpha=\beta=1$ case: prior proof \cite[Theorem~1]{karppolynomial}; other strict cases not covered.\\
Conjecture 2, p.~31 (Hurwitz stability) & The geometric case is the zero polynomial; the nonzero, strictly log-concave case is not covered.\\
Conjecture 3, p.~40 (removing discrete shifts in Theorems 5--6) & Not covered: its parameter shifts differ from the opposite degree shifts here.\\
Pr\'ekopa--Ninh comparison, pp.~52--54 & Not covered: it involves ${}_2F_2$ and shifts its parameters together.\\
General parameter classification, pp.~14--17; positive-valued restrictions & Not covered: no global sign or zero-free classification is given here.\\
\bottomrule
\end{tabular}
\end{center}

\subsection*{Verification scope}
\begingroup
\raggedright
All formal names below have prefix
\path{D5.S3.Analytic.GaussZeroBalanced}.
In the pinned \href{https://github.com/the-omega-institute/trureturing/blob/b8d708b1b89112061f30f218f031c1bf3672ebd8/D5/S3/Analytic/GaussZeroBalanced/FractionalLegendreTuran.lean}{\texttt{FractionalLegendreTuran.lean}},
the declaration \texttt{FractionalLegendreTuran.result}
proves $\neg\,\texttt{claim}$, where \texttt{claim} is precisely the
fully quantified strict comparison for $\mu>0$, $0<\delta<1$,
and $0<y<1$ \cite{lean}.
Its auxiliary declarations \texttt{minus\_enclosure},
\texttt{center\_enclosure}, and \texttt{plus\_enclosure} establish
the three intervals of Theorem~\ref{thm:main}; \texttt{witness\_negative}
gives strict negativity. These auxiliary declarations are private to
that file. The pinned
\href{https://github.com/the-omega-institute/trureturing/blob/b8d708b1b89112061f30f218f031c1bf3672ebd8/D5/S3/Analytic/GaussZeroBalanced/Defs.lean}{\texttt{Defs.lean}}
defines \texttt{Defs.gauss} as the real Gauss hypergeometric function;
the pinned
\href{https://github.com/the-omega-institute/trureturing/blob/b8d708b1b89112061f30f218f031c1bf3672ebd8/D5/S3/Analytic/GaussZeroBalanced/Series.lean}{\texttt{Series.lean}}
proves \texttt{Series.coefficient\_recurrence} for $c>0$.
The non-strict consequence and Corollaries~\ref{cor:open}--\ref{cor:zeros}
are written proofs. The integer recurrence evaluations and terminating
polynomial controls in \texttt{check\_bounds.py} are finite exact checks;
they do not prove the endpoint expansions or the noninteger-degree family.
\par\endgroup

\begingroup
\raggedright
\begin{thebibliography}{9}
\bibitem{karp} D.~Karp,
\emph{Log-convexity and log-concavity for series in product ratios of
rising factorials and gamma functions}, Special Functions in the 21st
Century: Theory and Applications, Washington, DC, April 6--8, 2011.
Conjecture~4, physical PDF pp.~52--54.
\url{https://math.nist.gov/~DLozier/SF21/SF21slides/Karp.pdf}.
\bibitem{dlmflegendre} NIST Digital Library of Mathematical Functions,
\S14.3(i), hypergeometric representation of the Ferrers function.
\url{https://dlmf.nist.gov/14.3.E1}.
\bibitem{dlmfhypergeometric} NIST Digital Library of Mathematical Functions,
\S15.2(i), Gauss series.
\url{https://dlmf.nist.gov/15.2.E1}.
\bibitem{karppolynomial} D.~Karp,
\emph{Positivity of Toeplitz determinants formed by rising factorial
series and properties of related polynomials}, arXiv:1203.1482v1,
March 7, 2012. Lemma~3 and Theorem~1.
\url{https://arxiv.org/abs/1203.1482v1}.
\bibitem{dlmfcontinuation} NIST Digital Library of Mathematical Functions,
Eq.~15.8.10, zero-balanced continuation formula ($m=0$).
\url{https://dlmf.nist.gov/15.8.E10}.
\bibitem{dlmfgamma} NIST Digital Library of Mathematical Functions,
Eqs.~5.5.2--5.5.3, digamma recurrence and gamma reflection.
\url{https://dlmf.nist.gov/5.5.E2};
\url{https://dlmf.nist.gov/5.5.E3}.
\bibitem{lean} \emph{FractionalLegendreTuran.lean}, declaration
\texttt{result}, fixed-revision source.
\url{https://github.com/the-omega-institute/trureturing/blob/b8d708b1b89112061f30f218f031c1bf3672ebd8/D5/S3/Analytic/GaussZeroBalanced/FractionalLegendreTuran.lean}.
\end{thebibliography}
\endgroup
\smallskip
\noindent\small\textcopyright\ 2026 Haobo Ma and Wenlin Zhang.
Licensed under Creative Commons Attribution 4.0 International (CC BY 4.0).
\end{document}
